% Generated draft: requires mathematical and source review.
\documentclass[11pt]{article}
\usepackage[margin=25mm]{geometry}
\usepackage{fontspec}
\setmainfont{Latin Modern Roman}
\usepackage{amsmath,amssymb,mathtools,unicode-math}
\setmathfont{Latin Modern Math}
\usepackage{enumitem,needspace,xurl,hyperref}
\setlength{\parindent}{0pt}
\setlength{\parskip}{5pt}
\setlength{\emergencystretch}{3em}
\newcommand{\N}{\mathbb{N}}
\newcommand{\Z}{\mathbb{Z}}
\newcommand{\Q}{\mathbb{Q}}
\newcommand{\R}{\mathbb{R}}
\newcommand{\C}{\mathbb{C}}
\newcommand{\F}{\mathbb{F}}
\newcommand{\A}{\mathbb{A}}
\newcommand{\PP}{\mathbb{P}}
\newcommand{\E}{\mathbb{E}}
\newcommand{\eps}{\varepsilon}
\newcommand{\dd}{\,\mathrm{d}}
\begin{document}
\section*{UnsolvedMath ID: 30002433}
\textbf{Source problem code:} catalogue entry
\subsection{Definitions and mathematical statement}
Let \(a,b\in\Z_{>0}\) satisfy \(\gcd(a,b)=1\), and put \(c=a+b\). The supplied passage proposes the following heuristic assumption for an unspecified refinement of the \(abc\) conjecture: the three kernels \(k(a),k(b),k(c)\) may be modeled as independent random variables.

The supplied record does not define \(k\). If ``kernel'' denotes the squarefree kernel, its definition would be
\[
 k(n)=\operatorname{rad}(n)=\prod_{\substack{p\text{ prime}\\p\mid n}}p
 \qquad(n\in\Z_{>0}),\qquad k(1)=1.
\]
This interpretation is conditional and requires source confirmation.

For fixed \(a,b\), these kernels are deterministic quantities. To formulate probabilistic independence meaningfully, one must specify a probability space, a sampling rule for coprime pairs, and the associated random variables. None is supplied; nor is any limiting regime, approximation criterion, or resulting refined \(abc\) bound stated. Consequently, the passage records a heuristic modeling assumption, not a precisely specified independence conjecture or mathematical problem. The supplied verification explicitly identifies it as a heuristic used to derive a proposed refinement.
\subsection{Short English statement}
For coprime positive integers \(a,b\), the source suggests treating the kernels of \(a\), \(b\), and \(a+b\) as independent random variables in a heuristic refinement of the \(abc\) conjecture. The kernel definition, probability model, and intended refined bound are missing, so a precise conjecture cannot be recovered from this record.
\subsection{Sources}
\begin{itemize}
\item Supplied local UnsolvedMath registry record, permanent ID \(30002433\), title ``Independent-Kernel Heuristic for Refined abc Bounds,'' source code \texttt{catalogue entry}, in the Oberwolfach Reports Open Problems collection. No individual proposer is supplied; the proposed year is recorded as 2013.
\item Supplied report citation: \emph{Analytic Number Theory (2014)}, Oberwolfach Reports, pp.~2963--3037, DOI \texttt{10.4171/owr/2013/51}. Source URL and supplied literature reference labeled ``original OWR report'': \url{https://doi.org/10.4171/owr/2013/51}.
\item Literature assessment dated \texttt{2026-08-21}: supplied dated curation, not independently verified here. It labels the record open and describes the passage as an unfinished heuristic refinement of the \(abc\) conjecture, rather than a verified theorem, with an unresolved underlying claim. The supplied background requests independent verification of time-sensitive status and cited claims.
\item The supplied statement-verification field reports comparison with the official report and says that independence is explicitly a heuristic assumption, not a separate mathematical question or conjecture. It promotes no clean statement and retains the original record for provenance; the supplied statement status is \texttt{unrecoverable}. These are attributed curation statements, not verification performed for this draft.
\end{itemize}
\end{document}
